% ECONOMICS_PROBLEM: {"id": "C73-15", "title": "Payoff characterization of convergence for a fixed mirror-descent rule", "jel_code": "C73", "source_id": "OP-0148"}
% FIRST_POSTED: 2026-10-11
% ATTEMPT_STATUS: solved
% ENTRY_KIND: research_attempt
% SUBMISSION: {"kind": "research_attempt", "category": "economics", "problem_id": "C73-15", "model": {"provider": "Anthropic", "version": "Claude Fable 5.1 (claude-fable-5-1)", "reasoning": "Claude Code agent session, default effort", "generated_on": "2026-10-11"}, "other_models": "None for the Lean development; the hand proof that is formalized was generated by GPT Codex (pull request #140, manuscript linked in the text); the game-model files are shared with the C73-11 formalization", "tools": "Lean 4.34.1 with Mathlib d13f23b723b8a846827a245b89c10fc7d3f11612 and lake; compilation on a separate Linux server", "target": "The exact C73-15 target as stated by the hand proof: for the fixed mirror-descent rule (arbitrary proper lower-semicontinuous strictly convex regularizers finite and differentiable on the relative interiors, arbitrary positive steps), a payoff-array characterization of universal convergence to the Nash set from every interior initialization, with all quantifiers retained", "scope": "Theorem 1 (assertions 1, 2, 3 equivalent) is formalized in full for every canonical rule and every payoff array for which the rule is well defined and gradient-domain preserving, as the manuscript assumes; the gradient enters only through its directional property on interior rays, as in the manuscript; the boundary test of Section 7 is not formalized", "human_contributions": "None; the submitter specified the task and submitted the result, with no mathematical edits", "verification": "Lean kernel check of every statement with only the standard axioms (propext, Classical.choice, Quot.sound); model self-audit of the definitions against the manuscript; no independent human proof verification", "reproducibility": "Public Lean repository pinned to a commit, with build commands and the actual checking outcome reported; the interactive agent-session prompts are not published"}
% FULL_SOLUTION_SCOPE: {"target": "The entire original C73-15 target as formulated by the hand proof (Theorem 1 of the manuscript)", "result": "main_theorem in Mirror/Main.lean: Converges R U <-> Criterion R U and Criterion R U <-> CriterionO R U, for a canonical rule R with R.WellDefined U; criterion_iff_no_chain gives the literal no-infinite-branch form", "coverage": "Every finite player set and nonempty finite action sets, every canonical rule (proper lsc strictly convex regularizers finite on the relative interiors, declared gradients with the directional property on interior rays, positive steps, any dense interior probe enumeration), every payoff array for which the rule is well defined and gradient-domain preserving from interior initializations; all interior initializations; the exact Nash set with Euclidean distance; both directions; no existence theorem for Nash equilibria is used", "dependencies": "Mathlib only (d13f23b7): continuity of convex functions on the intrinsic interior, compactness of countable products, well-founded ranks and the regularity of aleph_1; no additional axioms", "human_verification": "None", "unverified": "The transcription of the manuscript's definitions into Lean was checked by the model only; the hand proof itself is an LLM claim pending community expert review; the Section 7 boundary test is not formalized", "novelty": "Not checked beyond the manuscript's own references (Bichler, Legacci, Mertikopoulos, Oberlechner and Pradelski 2025)"}
\documentclass[11pt]{article}
\usepackage[margin=1in]{geometry}
\usepackage{amsmath,amssymb,amsthm}
\usepackage[hidelinks]{hyperref}
\newtheorem{theorem}{Theorem}
\theoremstyle{definition}
\newtheorem{definition}[theorem]{Definition}
\newcommand{\N}{\mathbb N}
\newcommand{\R}{\mathbb R}
\DeclareMathOperator{\NE}{NE}
\DeclareMathOperator{\dist}{dist}
\DeclareMathOperator{\ri}{ri}
\title{A complete machine-checked proof of the payoff characterization of universal convergence of fixed mirror descent}
\author{Claude Fable 5.1 (Claude Code), submitted by Ji Ho Bae}
\date{October 11, 2026}
\begin{document}
\maketitle
\begin{abstract}
The hand proof for this problem characterizes, for a fixed mirror-descent rule with arbitrary proper lower-semicontinuous strictly convex regularizers and arbitrary positive steps, the payoff arrays for which the exact simultaneous updates approach the Nash set from every interior initialization: by the well-foundedness of countable trees whose nodes are finite closed feasibility problems built from regularizer evaluations at interior probes, interior-ray secant certificates, the specified steps and explicit payoff polynomials. This note reports a complete Lean~4 formalization of that proof against Mathlib: the three assertions of the theorem are proved equivalent for every canonical rule and every payoff array for which the rule is well defined, as the manuscript assumes. Convergence to the Nash set is formulated literally, so that no existence theorem for Nash equilibria is needed. The development has no \texttt{sorry}, no warnings, assumes nothing beyond Mathlib, and uses only the standard axioms.
\end{abstract}

\section{The statement}

Fix players $N$ and nonempty finite action sets $S_i$, $X_i=\Delta(S_i)$, $X=\prod_iX_i$. The fixed rule has regularizers $h_i:X_i\to\R\cup\{+\infty\}$, proper, lower semicontinuous and strictly convex, finite and differentiable on $\ri X_i$, with a declared gradient $\nabla h_i$ on a declared domain containing $\ri X_i$, and steps $\eta_t>0$. For a payoff array $U$ with multilinear extensions $u_i^U$ and $v_i^U(x)_a=u_i^U(a,x_{-i})$, the exact simultaneous update is
\[
x_i^{t+1}=\arg\max_{p\in X_i}\bigl\{\eta_t\langle v_i^U(x^t),p\rangle-h_i(p)+\langle\nabla h_i(x_i^t),p\rangle\bigr\},
\]
under the canonical assumption that these updates are well defined and remain in the declared gradient domain from every interior initialization. For each $i$ fix an enumeration $(q_{i,\ell})$ of interior probes dense in $\ri X_i$ (the rational interior points). The trees $T_{a,b}(U)$, $a,b\ge1$, have nodes with finite discrete labels (strictly increasing marked times $T_j$, witnesses $(i_j,s_j)$, positive integer bounds $B_t$, positive integer denominators $N_{i,t,\ell,k}$); a node is valid when the finite system (f1)--(f5) of the manuscript has a solution in the variables $x^t\in X$, $c_{i,t}\in\R$, $g_{i,t}\in\R^{S_i}$: $x^1\in K_a=\{x:x_{i,s}\ge1/a\}$ and the bounds (f1); the epigraph and probe inequalities (f2); the interior-ray secant bounds (f3) with the labelled denominators; the update probe inequalities (f4) with score $g_{i,t}+\eta_tv_i^U(x^t)$; and the marked regret gaps $R^U_{i_j,s_j}(x^{T_j})\ge1/b$ (f5). A child appends one mark, retaining all earlier labels.

\begin{theorem}\label{main}
For every payoff array $U$ for which the canonical rule is well defined and gradient-domain preserving, the following are equivalent:
\textup{(1)} for every $x^1\in\ri X$ the trajectory of the rule satisfies $\dist(x^t,\NE(U))\to0$;
\textup{(2)} $T_{a,b}(U)$ is well founded for every $a,b\ge1$;
\textup{(3)} every $T_{a,b}(U)$ admits a strictly decreasing countable-ordinal ranking.
\end{theorem}

\section{What is verified}

The Lean development is public at a fixed commit \cite{Lean}. The rule is \texttt{Rule} (regularizer values on the effective domain, declared gradients and domain, steps, probe enumerations) and its standing assumptions are \texttt{Rule.Canonical}: proper lower semicontinuity and strict convexity of each $h_i$ on $X_i$, finiteness on $\ri X_i$, $\ri X_i\subseteq\mathrm{dom}_i$, and the directional property (dp) of the declared gradient along every interior ray, which is all the manuscript uses of differentiability. The update is \texttt{Rule.IsStep}, a trajectory is \texttt{Rule.Traj}, and the canonical assumption for $U$ is \texttt{Rule.WellDefined}; uniqueness of the update is proved from strict convexity (\texttt{Rule.IsArgmax.unique}) rather than assumed. Nodes, validity and the trees are \texttt{Node}, \texttt{Feasible}, \texttt{Valid}, \texttt{TreeWF} and \texttt{CondO}. Convergence \texttt{ConvNE} is the literal statement ``for every $\varepsilon>0$, eventually some Nash profile is within $\varepsilon$''.

\subsection*{The theorem}
\begin{verbatim}
theorem main_theorem (hR : R.Canonical) (hU : R.WellDefined U) :
    (Converges R U <-> Criterion R U) /\ (Criterion R U <-> CriterionO R U)
\end{verbatim}
Unicode is transliterated here (\texttt{<->}, \texttt{/\textbackslash}, \texttt{forall}, \texttt{->}). \texttt{Converges R U} is assertion (1) (every trajectory from every interior profile converges), \texttt{Criterion R U} is (2) (\texttt{TreeWF R U a b} for all $a,b\ge1$) and \texttt{CriterionO R U} is (3). The direction (2)$\Rightarrow$(1), \texttt{converges\_of\_criterion}, does not use the well-definedness assumption; (1)$\Rightarrow$(2), \texttt{criterion\_of\_converges}, uses it to identify the limit sequence with the declared rule. \texttt{criterion\_iff\_no\_chain} is (2) in the literal ``no infinite branch'' form.

\subsection*{The two lemmas of Section 2}
\begin{verbatim}
theorem Rule.certificate (hR : R.Canonical) (i : N) {x} (hx : x in R.fin i) {c : R} {g}
    (hc : R.h i x <= c) (hE : forall l, 0 <= R.h i (q l) - c - ip g (q l - x))
    (hN : forall l k, exists n, 1 <= n /\ 0 <= ... /\ ... <= enorm (q l - x) / ((k + 1) * n)) :
    c = R.h i x /\ forall l, Tendsto (fun t => (R.h i (ray x (q l) t) - R.h i x) / t)
      (nhdsWithin 0 (Ioi 0)) (nhds (ip g (q l - x)))

theorem Rule.isArgmax_of_probes (hR : R.Canonical) (i : N) {p y} (hp : p in R.fin i)
    (hprobe : forall l, 0 <= R.h i (q l) - R.h i p - ip y (q l - p)) : R.IsArgmax i y p
\end{verbatim}
(abbreviated, with \texttt{in} for membership; $q_\ell$ is \texttt{(R.probe i).q l}). The first is Lemma~1 (dense interior-ray certificate), the second Lemma~2 (dense-probe optimizer certificate); \texttt{Rule.ip\_eq\_of\_probes} identifies two vectors agreeing on all probe directions on every tangent direction of the simplex.

\section{Architecture of the formal proof}

\begin{itemize}
\item \emph{Convex analysis} (Sections 1--2). Continuity of $h_i$ on $\ri X_i$ is Mathlib's continuity of convex functions on the intrinsic interior (\texttt{ConvexOn.continuousOn\_intrinsicInterior}), after showing that $\ri X_i$ lies in the intrinsic interior of the effective domain. Lower semicontinuity makes the epigraph condition closed (\texttt{Rule.isClosed\_epi}). Secant quotients along rays into the effective domain are nondecreasing (\texttt{ConvexOn.secant\_mono}); (dp) therefore gives the supporting inequalities and integer secant denominators of any precision. Lemma~1 extends the probe inequalities to all interior points by density and continuity, then to the effective domain along segments towards an interior probe, and concludes with secant monotonicity; Lemma~2 is the same extension.
\item \emph{Nodes} (Section 4). Labels outside their ranges are normalized to zero, so the normalized nodes are determined by finite lists and the vertex set is countable (\texttt{Node.encode\_injective}); the constraints of a parent are among those of a child (\texttt{Feasible.of\_child}).
\item \emph{Failure produces a branch} (Section 5). Lemma~2 of the C73-11 development gives $b$ and the marked times; $c_{i,t}=h_i(x^t_i)$, the normalized declared gradient, integer bounds and the denominators from (dp) make every marked prefix a valid node (\texttt{mkNode\_feasible}).
\item \emph{A branch produces a trajectory} (Section 5). Witnesses of the nodes are placed in compact boxes outside their ranges; the product of the boxes over all dates is compact and first countable, so Mathlib's \texttt{IsCompact.tendsto\_subseq} yields the diagonal subsequence (\texttt{Chain.exists\_limit}). Each constraint is closed (for (f3) on the set where the profile is mixed, by \texttt{ContinuousOn.preimage\_isClosed\_of\_isClosed}) and holds along the subsequence from the node's depth on, hence in the limit. Lemma~1 and Lemma~2 at the limit, uniqueness of the maximizer, and the canonical assumption identify the limit sequence with the trajectory of the declared rule by induction (\texttt{not\_convNE\_of\_chain}); the marked inequalities survive.
\item \emph{Ordinals.} Well-foundedness of a relation on a countable type is equivalent to a strictly decreasing labeling by a countable ordinal (\texttt{WellFounded.rank}, regularity of $\aleph_1$).
\end{itemize}

\section{Scope}

The formalization follows the hand proof. The regularizer is represented by its finite values on the effective domain together with that domain, the manuscript's $h_i$ being the extension by $+\infty$. The gradient enters only through its directional property on interior rays, exactly as the manuscript stipulates, so Fr\'echet and finite-domain G\^ateaux gradients both qualify. The probe enumeration is any fixed dense sequence of interior points, of which the rational points are one (\texttt{exists\_probes}); probe indices and precisions are zero-based. Convergence to the Nash set is formulated as above, so no existence theorem for Nash equilibria is needed (the hand proof invokes Nash's theorem only for the distance lemma); when the Nash set is nonempty this is the convergence of Mathlib's \texttt{Metric.infDist} (\texttt{convNE\_iff\_infDist}). The boundary test of Section~7 of the hand proof (a specific regularizer) is not formalized. The formal theorems provide no finite recognition algorithm, and the hand proof does not claim one.

\section{Build and audit}

Lean \texttt{v4.34.1}, Mathlib commit \texttt{d13f23b723b8a846827a245b89c10fc7d3f11612}; about 2,550 lines in 11 files. The commands \texttt{lake exe cache get}, \texttt{lake build} and \texttt{lake env lean Verify.lean} were run on a fresh clone on a separate Linux machine. The build reports no errors and no warnings; the sources contain no \texttt{sorry}. Each of the 25 statements audited in \texttt{Verify.lean}, including \texttt{main\_theorem}, depends only on \texttt{propext}, \texttt{Classical.choice} and \texttt{Quot.sound}.

\section{Completion Estimate}
\noindent\textbf{COMPLETION: 100\%}

This is the claim for the full catalogue question as stated by the hand proof: an exact (infinitary) payoff-array characterization of universal convergence of the fixed mirror-descent rule to the Nash set. The machine check covers the three equivalences of Theorem~\ref{main}, for every canonical rule and every payoff array for which the rule is well defined, with no external theorem. It is not an independent expert review, and the hand proof itself remains an LLM claim pending community review.

\section*{Attribution and reproduction}
Model and process: Claude Fable 5.1 (Claude Code agent session, Anthropic), 11 October 2026, working from the manuscript \cite{Hand} and the Mathlib library; no other model contributed to the Lean development. Human contributions: none; the submitter specified the task and submitted the result, with no mathematical edits. The complete generated output is the Lean repository at its pinned commit \cite{Lean}; the interactive session prompts are not published. Lean source: \url{https://github.com/jbaelaw/mirror-descent-characterization-lean-20261011/tree/9c04e33b32ddbf8a072b04ebc850845040e5cde2} (\texttt{Mirror/Main.lean}, theorem \texttt{main\_theorem}). Reproduction: Lean \texttt{v4.34.1}, Mathlib commit \texttt{d13f23b723b8a846827a245b89c10fc7d3f11612}, built with \texttt{lake} as described above; the project has several files and depends on Mathlib, so it is not a single-file Lean Web example. Checking outcome: no errors, no warnings, no \texttt{sorry}, standard axioms only, on a fresh clone. The repository maintainers have not verified the result.

\section*{Final status}
SOLVED, as a machine-checked formalization of the hand proof's complete characterization; community expert review of the characterization is pending.

\begin{thebibliography}{9}
\bibitem{Hand} Ji Ho Bae, \emph{A closed-certificate characterization of universal convergence of fixed mirror descent}, hand proof for C73-15, \url{https://github.com/jbaelaw/fixed-mirror-convergence-characterization/blob/000960f40cf5f7565032a9ba5a2986c93b5e0c5b/paper/Proof.tex}.
\bibitem{Lean} Lean 4 formalization, \url{https://github.com/jbaelaw/mirror-descent-characterization-lean-20261011/tree/9c04e33b32ddbf8a072b04ebc850845040e5cde2}.
\bibitem{BLMOP} Martin Bichler, Davide Legacci, Panayotis Mertikopoulos, Matthias Oberlechner and Bary Pradelski, \emph{Characterizing the convergence of game dynamics via potentialness}, arXiv:2503.16285 (2025).
\end{thebibliography}
\end{document}
